\documentclass[a4paper,12pt,portuguese]{report}

\usepackage[T1]{fontenc}
\usepackage[scaled]{luximono}
\usepackage{amsmath}
\usepackage{amssymb}
\usepackage{babel}
\usepackage[applemac]{inputenc}
\usepackage{graphicx,color}
\usepackage{epstopdf}
\DeclareGraphicsRule{.tif}{png}{.png}{`convert #1 `dirname #1`/`basename #1 .tif`.png}

\usepackage{marvosym}
\usepackage[Lenny]{fncychap}
\usepackage{listings}
\lstloadlanguages{Python,C,Java,Prolog,Lisp}
\usepackage{pifont}
\usepackage[colorlinks,urlcolor=blue]{hyperref}

\usepackage{verbatim}
\usepackage{alltt}
\usepackage[official]{eurosym}

\pdfpagewidth=\paperwidth
\pdfpageheight=\paperheight

\newcommand{\python}{\textbf {Python}}

\lstset{frame=single,language=Python,captionpos=b,showstringspaces=false,basicstyle=\ttfamily,commentstyle=\color{blue},tabsize=2,breaklines=true, extendedchars=true}
%----------------------------------------------

\definecolor{cinza}{rgb}{0.192,0.192,0.192}

\definecolor{orange}{rgb}{1,0.5,0.0}


\definecolor{rose}{rgb}{1,0.6,0.6}

\newcommand{\problema}[1]{\noindent \fbox{} \colorbox{cinza}{\textbf{\textcolor[rgb]{1.00,1.00,1.00}{Problema #1}}}}


\newcommand{\modulo}[1]{\noindent\colorbox{cinza}{\textbf{\textcolor[rgb]{1.00,1.00,1.00}{M�dulo #1}}}}

\newcommand{\muitofacil}{\noindent \colorbox{magenta}{\textbf{\textcolor[rgb]{1.00,1.00,1.00}{MF}}}}

\newcommand{\facil}{\noindent \colorbox{green}{\textbf{\textcolor[rgb]{1.00,1.00,1.00}{F}}}}

\newcommand{\medio}{\noindent \colorbox{orange}{\textbf{\textcolor[rgb]{1.00,1.00,1.00}{M}}}}

\newcommand{\dificil}{\noindent \colorbox{red}{\textbf{\textcolor[rgb]{1.00,1.00,1.00}{D}}}}

\newcommand{\muitodificil}{\noindent \colorbox{rose}{\textbf{\textcolor[rgb]{1.00,1.00,1.00}{MD}}}}

%-----------------------------------------------
\renewcommand\lstlistingname{Listagem}

\lstset{backgroundcolor=\color{yellow},language=Python,captionpos=b,showstringspaces=false,basicstyle=\ttfamily,commentstyle=\color{blue},tabsize=2,breaklines=true, extendedchars=true}


%\lstset{language=Python,backgroundcolor=\color{yellow},captionpos=b}
% ----------------------------------------------------------------
\newenvironment{objectivos}
    {\noindent {\bfseries{\Large{Objectivos}}}
        \begin{dinglist}{"34}}
    {\end{dinglist}}

% contador de Exercicios
\newcounter{exercicio}[chapter]
\renewcommand\theexercicio{\thechapter.\arabic{exercicio}}


\newenvironment{Exercicio}
{\noindent  \refstepcounter{exercicio}  \colorbox{blue}{  \textbf{\textcolor[rgb]{1.00,1.00,1.00}
{ Problema \theexercicio}}}}
{ \vspace*{0.2cm} \par}

\newenvironment{requisitos}
    {\noindent {\bfseries{\Large{Requisitos}}}
        \begin{dinglist}{"34}}
    {\end{dinglist}}

\newenvironment{materialapoio}
    {\noindent {\bfseries{\Large{Material de Apoio}}}
        \begin{dinglist}{"34}}
    {\end{dinglist}}

\newenvironment{sol}
{\noindent  \colorbox{magenta}{\textbf{\textcolor[rgb]{1.00,1.00,1.00}
{ Solu��o}}}}
{ \vspace*{0.2cm} \par}




\begin{document}

\noindent \rule{\linewidth}{1mm}

\begin{center}

\begin{figure}[htbp]
 \centering
 \includegraphics{timbre.eps}
 \end{figure}

\LARGE{\textbf{\textcolor[rgb]{0.98,0.00,0.00}
{Introdu��o � Programa��o \\ e Resolu��o de Problemas}}} \\
\vspace*{0.25cm}2009/2010 \\\vspace*{0.25cm}

\rule{\linewidth}{1mm}

\vspace*{5.0cm} \textcolor[rgb]{0.00,0.40,0.29}{Ficha \#6 }

\vspace*{2.0cm}
\noindent Ernesto Costa,\\ Lu�s Macedo, Rog�rio de Lemos, Vasco Pereira
\end{center}

\renewcommand\chaptername{Ficha}
\setcounter{chapter}{9}
\chapter{Exerc\'{i}os V\'{a}rios}
\setcounter{chapter}{6}



\begin{objectivos}
\item
\end{objectivos}

\bigskip

\begin{requisitos}
\item

\end{requisitos}
\bigskip

\begin{materialapoio}
\item

\end{materialapoio}




\newpage
\section{O Jogo}

O Mundo do DEI \'{e} um jogo de aventura muito simples
baseado em texto, onde os utilizadores podem perambular 
nas diferentes espa\c{c}os de um edif\'{i}cio. Este jogo deve ser estendido para se tornar mais interessante.

Nesta vers\~{a}o simples do Mundo do DEI, o DEI \'{e} divido em salas, e o jogador tem que descobrir qual a configura\c{c}\~{a}o do DEI percorrendo as diversas salas. Estas est\~{a}o interconectadas de diferentes maneiras atrav\'{e}s de portas. Nesta vers\~{a}o do jogo \'{e} assumido que o DEI tem quatro salas, excluindo a parte de fora do edif\'{i}cio - onde sempre se come\c{c}a o jogo. Os tr\^{e} comandos disponiveis s\~{a}o os seguintes:

\begin{itemize}
	\item \emph{go} - sai de uma da sala para outra seguindo uma determinada dire\c{c}\~{a}o que inicialmente \'{e} uma destas quatro: \emph{norte}, \emph{sul}, \emph{leste} e \emph{oeste}.
\item \emph{quit} - permite sair do jogo.
\item \emph{help} -  fornece informa\c{c}\~{a}o muito elementar sobre os diversos comandos 
\end{itemize}

Antes de fazer os seguintes exerc\'{i}cio que implecam na modifica\c{c}\~{a}o do c\'{o}digo, aconselha-se que joguem o jogo e entedam muito bem o c\'{o}digo \textbf{gameDEI.py} que est\'{a} dispon\'{i}vel no \textbf{WoC}, na p\'{a}gina da cadeira, e que devem ser descarregado por si.


\begin{Exercicio}

Na primeira modifica\c{c}\~{a}o, considere que uma porta \'{e}' aberta no edif\'{i}cio do DEI que vai do exterior do edif\'{i}cio diretamente \`{a} biblioteca. Fa\c{c}a esta adi\c{c}\~{a}o ao programa, e considere que a porta est\'{a} ao \emph{sudoeste} do edif\'{i}cio.
\end{Exercicio}

\begin{Exercicio}

Para a segunda modifica\c{c}\~{a}o, considere que uma sala adicional \'{e} acrecentada ao DEI. Esta seria uma sala de jogos que est\'{a} ligada pelo sul ao corredor, pelo norte ao caf\'{e}, e pelo leste \`{a} biblioteca. Fa\c{c}a a devida modifica\c{c}\~{a}o no programa.
\end{Exercicio}

\begin{Exercicio}

Uma das defici\^{e}ncias deste programa \'{e} que somente aceita comandos em letras min\'{u}sculas. Fa\c{c}a as devidas modifica\c{c}\~{o}es para permitir o uso tanto de min\'{u}sculas como de mai\'{u}sculas. 

\end{Exercicio}

\begin{Exercicio}

Espalhadas pelo edif\'{i}cio existem tr\^{e}s chaves (\emph{ouro}, \emph{prata} e \emph{bronze}). Uma vez que o jogador esteja numa sala que contenha uma das chaves, o jogador pode pegar a chave dispon\'{i}vel. O jogador ganha o jogo se conseguir sair do edif\'{i}cio com as tr\^{e}s chaves. Para a implementa\c{c}\~{a}o desta modifica\c{c}\~{a}o sera necess\'{a}rio introduzir o comando "\emph{get}", que possibilita pegar uma chave. Para esta vers\~{a}o do jogo  \'{e} considerado que existe somente uma chave por sala (mas nem todas as salas t\^{e}m uma chave).

\end{Exercicio}

\begin{Exercicio}

Para se introduzir um componente surpresa no jogo, decide-se fazer a aloca\c{c}\~{a}o das chaves pelas salas de uma forma aleat\'{o}ria logo ao in\'{i}cio do jogo. Implemente esta modifica\c{c}\~{a}o no programa.

\end{Exercicio}

\begin{Exercicio}

Finalmente, para tornar o jogo mais emociante, estabelece-se que o verdadeiro ganhador s\'{o} pode demorar no m\'{a}ximo 2 minutos para finalizar a tarefa. Implemente esta modifica\c{c}\~{a}o no programa.

\end{Exercicio}

\begin{Exercicio}

Entretanto, muito ainda resta por fazer. Expanda o jogo como quiser: aumente o n\'{u}mero de salas, aumente o n\'{u}mero de chaves, possibilite jogar com mais de um jogador, monte uma interface gr\'{a}fica para o jogo, etc. 
\end{Exercicio}


\end{document} 